% ============================================================
%  SCHOLA DOMESTICA — Magister Mathematica
%  Level 1 | Week 11 | Day 4
%  Topic: Addition Facts to 10 — Addition Table & Word Problems
% ============================================================
\documentclass[12pt, letterpaper]{article}

\usepackage[top=1in, bottom=1in, left=1in, right=1in]{geometry}
\usepackage[T1]{fontenc}
\usepackage[utf8]{inputenc}
\usepackage{lmodern}
\usepackage{microtype}
\usepackage{amsmath}
\usepackage{xcolor}
\usepackage{tikz}
\usepackage{array}
\usepackage{enumitem}
\usepackage{multicol}
\usepackage{mdframed}
\usepackage{fancyhdr}

\definecolor{scholared}{RGB}{160,30,30}
\definecolor{scholagold}{RGB}{180,140,40}
\definecolor{scholacream}{RGB}{253,248,235}
\definecolor{scholagreen}{RGB}{50,110,60}
\definecolor{scholablue}{RGB}{40,80,140}

\pagestyle{fancy}
\fancyhf{}
\renewcommand{\headrulewidth}{0.6pt}
\fancyhead[L]{\small\textcolor{scholared}{\textbf{Schola Domestica — Magister Mathematica}}}
\fancyhead[R]{\small\textcolor{scholared}{Level 1 · Week 11 · Day 4}}
\fancyfoot[C]{\small\textcolor{gray}{— \thepage\ —}}

\newcommand{\answerline}{\underline{\hspace{2.5cm}}}
\newcommand{\biganswerline}{\underline{\hspace{4cm}}}
\newcommand{\numbox}{\tikz{\draw[scholablue,thick](0,0)rectangle(1.2cm,0.85cm);}}

\newmdenv[
  backgroundcolor=scholacream,
  linecolor=scholared,
  linewidth=1.5pt,
  roundcorner=6pt,
  innertopmargin=8pt,
  innerbottommargin=8pt,
  innerleftmargin=10pt,
  innerrightmargin=10pt
]{lessonbox}

\newmdenv[
  backgroundcolor={rgb,255:red,235;green,245;blue,235},
  linecolor=scholagreen,
  linewidth=1.2pt,
  roundcorner=4pt,
  innertopmargin=6pt,
  innerbottommargin=6pt,
  innerleftmargin=8pt,
  innerrightmargin=8pt
]{enrichbox}

\begin{document}

\begin{center}
  {\LARGE\textbf{\textcolor{scholared}{The Addition Table:}}}\\[4pt]
  {\Large\textcolor{scholagold}{\textit{Patterns in our Facts}}}\\[6pt]
  {\small Name: \biganswerline \hspace{1cm} Date: \answerline}
\end{center}

\vspace{4pt}
\noindent\textcolor{scholared}{\rule{\linewidth}{1pt}}
\vspace{4pt}

% IMAGE PROMPT (Miyazaki watercolour style):
% An elderly tanuki sits before a large stone tablet carved with numbers in neat rows, tracing
% them with one claw. Firefly light illuminates the carvings. Night forest, soft glowing palette,
% white paper background, no text.

\begin{lessonbox}
\textbf{An addition table} shows all the sums in one place.
Read across the top row to find the first addend.
Read down the left column to find the second addend.
Where they meet is the \textbf{sum}.\\[4pt]
\noindent The \textbf{diagonal} (top-left to bottom-right) shows the \emph{doubles}.
\end{lessonbox}

\bigskip

% ===== ADDITION TABLE (partial, 0–5 × 0–5) =====
\noindent\textcolor{scholared}{\large\textbf{Part I — Complete the Addition Table}} \hfill \textit{(20 cells)}

\medskip
\noindent\textit{Fill in every empty cell with the sum. The first row and column are done for you.}

\bigskip

\noindent
\renewcommand{\arraystretch}{1.6}
\begin{tabular}{|>{\columncolor{scholacream}}c|c|c|c|c|c|c|}
\hline
\rowcolor{scholacream}
\textbf{$+$} & \textbf{0} & \textbf{1} & \textbf{2} & \textbf{3} & \textbf{4} & \textbf{5} \\
\hline
\textbf{0} & 0 & 1 & 2 & 3 & 4 & 5 \\
\hline
\textbf{1} & 1 & \phantom{0}\underline{\hspace{0.6cm}} & \phantom{0}\underline{\hspace{0.6cm}} & \phantom{0}\underline{\hspace{0.6cm}} & \phantom{0}\underline{\hspace{0.6cm}} & \phantom{0}\underline{\hspace{0.6cm}} \\
\hline
\textbf{2} & 2 & \phantom{0}\underline{\hspace{0.6cm}} & \phantom{0}\underline{\hspace{0.6cm}} & \phantom{0}\underline{\hspace{0.6cm}} & \phantom{0}\underline{\hspace{0.6cm}} & \phantom{0}\underline{\hspace{0.6cm}} \\
\hline
\textbf{3} & 3 & \phantom{0}\underline{\hspace{0.6cm}} & \phantom{0}\underline{\hspace{0.6cm}} & \phantom{0}\underline{\hspace{0.6cm}} & \phantom{0}\underline{\hspace{0.6cm}} & \phantom{0}\underline{\hspace{0.6cm}} \\
\hline
\textbf{4} & 4 & \phantom{0}\underline{\hspace{0.6cm}} & \phantom{0}\underline{\hspace{0.6cm}} & \phantom{0}\underline{\hspace{0.6cm}} & \phantom{0}\underline{\hspace{0.6cm}} & \phantom{0}\underline{\hspace{0.6cm}} \\
\hline
\textbf{5} & 5 & \phantom{0}\underline{\hspace{0.6cm}} & \phantom{0}\underline{\hspace{0.6cm}} & \phantom{0}\underline{\hspace{0.6cm}} & \phantom{0}\underline{\hspace{0.6cm}} & \phantom{0}\underline{\hspace{0.6cm}} \\
\hline
\end{tabular}
\renewcommand{\arraystretch}{1}

\medskip
\noindent\textit{\small Colour the diagonal (doubles row) in yellow when you have finished.}

\bigskip\bigskip

% ===== PART II: USE THE TABLE =====
\noindent\textcolor{scholared}{\large\textbf{Part II — Use Your Table}} \hfill \textit{(Problems 1–8)}

\medskip
\noindent\textit{Find each fact in the table. Write the sum.}

\bigskip

\noindent
\begin{tabular}{@{}p{3.3cm} p{3.3cm} p{3.3cm} p{3.3cm}@{}}
\textbf{1.}\; $4 + 5 =$ \numbox &
\textbf{2.}\; $3 + 4 =$ \numbox &
\textbf{3.}\; $5 + 3 =$ \numbox &
\textbf{4.}\; $2 + 5 =$ \numbox \\[12pt]
\textbf{5.}\; $5 + 4 =$ \numbox &
\textbf{6.}\; $3 + 3 =$ \numbox &
\textbf{7.}\; $4 + 4 =$ \numbox &
\textbf{8.}\; $5 + 5 =$ \numbox
\end{tabular}

\bigskip\bigskip

% ===== PART III: WORD PROBLEMS =====
\noindent\textcolor{scholared}{\large\textbf{Part III — Story Problems}} \hfill \textit{(Problems 9–12)}

\medskip

\noindent\textbf{9.} Five children sit on one bench; four sit on another. How many children in all?

\medskip\noindent\hspace{2em} $\numbox + \numbox = \numbox$ children

\bigskip

\noindent\textbf{10.} A baker made \textbf{8} rolls and \textbf{2} loaves. How many baked goods in all?

\medskip\noindent\hspace{2em} $\numbox + \numbox = \numbox$ baked goods

\bigskip

\noindent\textbf{11.} Tom has \textbf{3} blue marbles and \textbf{7} red marbles. How many marbles altogether?

\medskip\noindent\hspace{2em} $\numbox + \numbox = \numbox$ marbles

\bigskip

\noindent\textbf{12.} Write your own problem with a sum of 10:

\medskip\noindent\hspace{2em}\underline{\hspace{9cm}}

\medskip\noindent\hspace{2em} $\numbox + \numbox = 10$

\bigskip

\begin{enrichbox}
\textbf{\textcolor{scholagreen}{Pattern: The table is symmetric!}}\\[3pt]
Look at $3 + 4$ and $4 + 3$ in your table. Are they the same? \quad$\square$\,Yes \; $\square$\,No\\[3pt]
\textit{This shows the turn-around rule (commutative property) at work for every fact!}
\end{enrichbox}

\vspace{0.3cm}
\noindent\textcolor{gray}{\small\textit{Ray's Primary Arithmetic — the complete addition table.}}

\end{document}
